% ============================================================================
% WORKBOOK — Section 3.2: Adding Integers
% CCSS 7.NS.A.1.B
% Practice-focused reinforcement for the study guide topic
% ============================================================================
\section{Adding Integers}
\topicTitle{Adding Integers}

\begin{quickReview}{Adding Integers}
\begin{itemize}[leftmargin=*]
    \item \textbf{Same signs:} Add the absolute values. Keep the common sign.
    
    $(-5) + (-3) = -8$ \qquad $4 + 7 = 11$
    \item \textbf{Different signs:} Subtract the smaller absolute value from the larger. Keep the sign of the number with the larger absolute value.
    
    $(-9) + 4 = -5$ \qquad $6 + (-2) = 4$
    \item \textbf{Additive inverses:} A number plus its opposite always equals $0$: \quad $n + (-n) = 0$.
\end{itemize}

\medskip
\textbf{Example:} $(-7) + 10 = ?$ \quad Different signs $\to$ $10 - 7 = 3$. Larger absolute value is $10$ (positive), so the answer is $3$.
\end{quickReview}

% ── Warm-Up ──────────────────────────────────────────────────────────────────
\begin{practiceBox}[funGreen]{\faSun~Warm-Up}

\practiceHeader[funGreen]{Add these integers.}

\begin{multicols}{2}
\prob $5 + 3$ \answerBlank[2cm]
\answerExplain{$8$}{Same signs (both positive), so add the absolute values: $5 + 3 = 8$. The result is positive.}

\prob $(-4) + (-6)$ \answerBlank[2cm]
\answerExplain{$-10$}{Same signs (both negative), so add: $4 + 6 = 10$. Keep the negative sign: $-10$.}

\prob $9 + (-9)$ \answerBlank[2cm]
\answerExplain{$0$}{These are additive inverses (opposites). A number plus its opposite always equals $0$.}

\prob $(-3) + 8$ \answerBlank[2cm]
\answerExplain{$5$}{Different signs, so subtract: $8 - 3 = 5$. The larger absolute value ($8$) is positive, so the answer is $5$.}

\prob $7 + (-10)$ \answerBlank[2cm]
\answerExplain{$-3$}{Different signs, so subtract: $10 - 7 = 3$. The larger absolute value ($10$) is negative, so the answer is $-3$.}

\prob $(-1) + (-1)$ \answerBlank[2cm]
\answerExplain{$-2$}{Same signs (both negative), so add: $1 + 1 = 2$. Keep the negative sign: $-2$.}
\end{multicols}
\end{practiceBox}

% ── Practice A: Same Signs ───────────────────────────────────────────────────
\begin{practiceBox}{Adding Integers with the Same Sign}

\practiceHeader{Add. Both numbers have the same sign.}

\begin{multicols}{2}
\prob $(-12) + (-8)$ \answerBlank[2cm]
\answerExplain{$-20$}{Same signs (both negative), so add absolute values: $12 + 8 = 20$. Keep the negative sign: $-20$.}

\prob $15 + 23$ \answerBlank[2cm]
\answerExplain{$38$}{Same signs (both positive), so add: $15 + 23 = 38$. The result is positive.}

\prob $(-7) + (-14)$ \answerBlank[2cm]
\answerExplain{$-21$}{Same signs (both negative), so add absolute values: $7 + 14 = 21$. Keep the negative sign: $-21$.}

\prob $(-25) + (-30)$ \answerBlank[2cm]
\answerExplain{$-55$}{Same signs (both negative), so add: $25 + 30 = 55$. Keep the negative sign: $-55$.}
\end{multicols}
\end{practiceBox}

% ── Practice B: Different Signs ──────────────────────────────────────────────
\begin{practiceBox}[funTeal]{Adding Integers with Different Signs}

\practiceHeader[funTeal]{Add. The numbers have different signs.}

\begin{multicols}{2}
\prob $(-15) + 9$ \answerBlank[2cm]
\answerExplain{$-6$}{Different signs, so subtract: $15 - 9 = 6$. The larger absolute value ($15$) is negative, so the answer is $-6$.}

\prob $20 + (-13)$ \answerBlank[2cm]
\answerExplain{$7$}{Different signs, so subtract: $20 - 13 = 7$. The larger absolute value ($20$) is positive, so the answer is $7$.}

\prob $(-50) + 35$ \answerBlank[2cm]
\answerExplain{$-15$}{Different signs, so subtract: $50 - 35 = 15$. The larger absolute value ($50$) is negative, so the answer is $-15$.}

\prob $11 + (-18)$ \answerBlank[2cm]
\answerExplain{$-7$}{Different signs, so subtract: $18 - 11 = 7$. The larger absolute value ($18$) is negative, so the answer is $-7$.}

\prob $(-6) + 6$ \answerBlank[2cm]
\answerExplain{$0$}{These are opposites (additive inverses), so their sum is $0$.}

\prob $(-100) + 45$ \answerBlank[2cm]
\answerExplain{$-55$}{Different signs, so subtract: $100 - 45 = 55$. The larger absolute value ($100$) is negative, so the answer is $-55$.}
\end{multicols}
\end{practiceBox}

% ── Practice C: Three or More Integers ───────────────────────────────────────
\begin{practiceBox}[funOrange]{Adding Three or More Integers}

\practiceHeader[funOrange]{Find each sum.}

\prob $(-4) + 7 + (-3)$ \answerBlank[2cm]
\answerExplain{$0$}{Add left to right: $(-4) + 7 = 3$. Then $3 + (-3) = 0$.}

\prob $5 + (-8) + (-2) + 10$ \answerBlank[2cm]
\answerExplain{$5$}{Group positives: $5 + 10 = 15$. Group negatives: $(-8) + (-2) = -10$. Combine: $15 + (-10) = 5$.}

\prob $(-11) + (-3) + 7 + (-1)$ \answerBlank[2cm]
\answerExplain{$-8$}{Group negatives: $(-11) + (-3) + (-1) = -15$. Add the positive: $-15 + 7 = -8$.}

\prob $6 + (-6) + (-9) + 9$ \answerBlank[2cm]
\answerExplain{$0$}{Notice the pairs of opposites: $6 + (-6) = 0$ and $(-9) + 9 = 0$. Total: $0 + 0 = 0$.}

\end{practiceBox}

% ── Practice D: True or False? ───────────────────────────────────────────────
\begin{practiceBox}[funPurple]{True or False?}

\trueOrFalse{The sum of two negative integers is always negative.}
\answerExplain{True}{When both integers are negative, you add their absolute values and keep the negative sign. The result is always negative.}

\trueOrFalse{The sum of a positive and a negative integer is always negative.}
\answerExplain{False}{The sign of the sum depends on which number has the larger absolute value. For example, $10 + (-3) = 7$ is positive.}

\trueOrFalse{$(-4) + 4 = 0$}
\answerExplain{True}{$-4$ and $4$ are opposites (additive inverses). A number plus its opposite always equals $0$.}

\end{practiceBox}

% ── Practice E: Word Problems ────────────────────────────────────────────────
\begin{practiceBox}[funBlue]{\faGlobe~Word Problems}

\wordProblem{The temperature was $-5°$C in the morning. By afternoon it rose $12°$C. What is the afternoon temperature?}{°C}
\answerExplain{$7°$C}{Add the rise to the starting temperature: $(-5) + 12 = 7$. Different signs, subtract: $12 - 5 = 7$, positive because $12$ has the larger absolute value.}

\wordProblem{A football team gained $8$ yards on the first play and lost $11$ yards on the second play. What was their total yardage?}{yards}
\answerExplain{$-3$ yards (a net loss of $3$ yards)}{Gain = $+8$, loss = $-11$. Add: $8 + (-11) = -3$. Different signs, subtract: $11 - 8 = 3$. Negative because $11$ has the larger absolute value.}

\wordProblem{An elevator starts at floor $-2$ (underground parking). It goes up $5$ floors, then down $3$ floors. What floor is it on now?}{floor}
\answerExplain{Floor $0$ (ground level)}{Start at $-2$. Go up $5$: $(-2) + 5 = 3$. Go down $3$: $3 + (-3) = 0$. The elevator is at ground level.}

\end{practiceBox}

% ── Challenge ────────────────────────────────────────────────────────────────
\begin{challengeBox}

\prob Find two integers whose sum is $-3$ and whose absolute values add up to $11$. \answerBlank[3cm]
\answerExplain{$4$ and $-7$ (or $-7$ and $4$)}{If $a + b = -3$ and $|a| + |b| = 11$, and the signs are different (since the sum is close to $0$), try $a = 4, b = -7$: $4 + (-7) = -3$ and $|4| + |{-7}| = 4 + 7 = 11$. \checkmark}

\prob The sum of three consecutive integers is $-12$. What are the three integers? \answerBlank[3cm]
\answerExplain{$-5, \; -4, \; -3$}{Consecutive integers: $n, n+1, n+2$. Their sum: $3n + 3 = -12$, so $3n = -15$, $n = -5$. The integers are $-5, -4, -3$. Check: $(-5) + (-4) + (-3) = -12$. \checkmark}

\end{challengeBox}

\encouragement{Great work adding integers! Positive or negative, you've got it covered!}
